\documentclass[a4paper,10pt,BCOR10mm,oneside,headsepline]{scrartcl}
\usepackage[portuges]{babel}
\usepackage[utf8]{inputenc}
\usepackage{amsmath}
\usepackage{wasysym}% provides \ocircle and \Box
\usepackage{enumitem}% easy control of topsep and leftmargin for lists
\usepackage{color}% used for background color
\usepackage{forloop}% used for \Qrating and \Qlines
\usepackage{ifthen}% used for \Qitem and \QItem
\usepackage{typearea}
\usepackage{graphicx}
\usepackage{booktabs}
\usepackage{listings}
\usepackage{color}
\usepackage{algorithm}
\usepackage{algorithmic}
\usepackage{amssymb}
\usepackage{minted}
\usepackage{mdframed}
\usepackage{wasysym}
\usepackage{xcolor}
\usepackage{fontspec}
%\usepackage{tgtermes}

\areaset{17cm}{26cm}
\setlength{\topmargin}{-2cm}

\definecolor{codegreen}{rgb}{0,0.6,0}
\definecolor{codegray}{rgb}{0.5,0.5,0.5}
\definecolor{codepurple}{rgb}{0.58,0,0.82}
\definecolor{backcolour}{rgb}{0.95,0.95,0.92}
 
\lstdefinestyle{mystyle}{
    backgroundcolor=\color{backcolour},   
    commentstyle=\color{codegreen},
    keywordstyle=\color{magenta},
    numberstyle=\tiny\color{codegray},
    stringstyle=\color{codepurple},
    basicstyle=\footnotesize,
    breakatwhitespace=false,         
    breaklines=true,                 
    captionpos=b,                    
    keepspaces=true,                 
    numbers=left,                    
    numbersep=5pt,                  
    showspaces=false,                
    showstringspaces=false,
    showtabs=false,                  
    tabsize=2
}
 
\lstset{style=mystyle}

\mdfsetup{
    skipabove=12pt,
    skipbelow=12pt,
    innertopmargin=8pt,
    innerbottommargin=8pt,
    innerleftmargin=10pt,
    innerrightmargin=10pt,
    splittopskip=12pt,
    splitbottomskip=12pt,
}

\algsetup{linenodelimiter=.,linenosize=\tiny}

\renewcommand{\algorithmicrequire}{\textbf{Entrada:}}
\renewcommand{\algorithmicensure}{\textbf{Saída:}}
\floatname{algorithm}{Algoritmo}

\setmainfont{TeX Gyre Termes}
\setmonofont{APL 385 Unicode}[Scale=0.9]

%% \Qq = Questionaire question. Oh, this is just too simple. It helps
%% making it easy to globally change the appearance of questions.
\newcommand{\Qq}[1]{\textbf{#1}}
%\newcommand{\Qq}[1]{{#1}}

%% \QO = Circle or box to be ticked. Used both by direct call and by
%% \Qrating and \Qlist.
\newcommand{\QO}{$\Box$}% or: $\ocircle$

%% \Qrating = Automatically create a rating scale with NUM steps, like
%% this: 0--0--0--0--0.
\newcounter{qr}
\newcommand{\Qrating}[1]{\QO\forloop{qr}{1}{\value{qr} < #1}{---\QO}}

%% \Qline = Again, this is very simple. It helps setting the line
%% thickness globally. Used both by direct call and by \Qlines.
\newcommand{\Qline}[1]{\noindent\rule{#1}{0.6pt}}

%% \Qlines = Insert NUM lines with width=\linewith. You can change the
%% \vskip value to adjust the spacing.
\newcounter{ql}
\newcommand{\Qlines}[1]{\forloop{ql}{0}{\value{ql}<#1}{\vskip0em\Qline{\linewidth}}}

%% \Qlist = This is an environment very similar to itemize but with
%% \QO in front of each list item. Useful for classical multiple
%% choice. Change leftmargin and topsep accourding to your taste.
\newenvironment{Qlist}{%
\renewcommand{\labelitemi}{\QO}
\begin{itemize}[leftmargin=1.5em,topsep=-.5em]
}{%
\end{itemize}
}

%% \Qtab = A "tabulator simulation". The first argument is the
%% distance from the left margin. The second argument is content which
%% is indented within the current row.
\newlength{\qt}
\newcommand{\Qtab}[2]{
\setlength{\qt}{\linewidth}
\addtolength{\qt}{-#1}
\hfill\parbox[t]{\qt}{\raggedright #2}
}

%% \Qitem = Item with automatic numbering. The first optional argument
%% can be used to create sub-items like 2a, 2b, 2c, ... The item
%% number is increased if the first argument is omitted or equals 'a'.
%% You will have to adjust this if you prefer a different numbering
%% scheme. Adjust topsep and leftmargin as needed.
\newcounter{itemnummer}
\newcommand{\Qitem}[2][]{% #1 optional, #2 notwendig
\ifthenelse{\equal{#1}{}}{\stepcounter{itemnummer}}{}
\ifthenelse{\equal{#1}{a}}{\stepcounter{itemnummer}}{}
\begin{enumerate}[topsep=2pt,leftmargin=2.8em]
\item[\textbf{\arabic{itemnummer}#1.}] #2
\end{enumerate}
}

%% \QItem = Like \Qitem but with alternating background color. This
%% might be error prone as I hard-coded some lengths (-5.25pt and
%% -3pt)! I do not yet understand why I need them.
\definecolor{bgodd}{rgb}{0.8,0.8,0.8}
\definecolor{bgeven}{rgb}{0.9,0.9,0.9}
\newcounter{itemoddeven}
\newlength{\gb}
\newcommand{\QItem}[2][]{% #1 optional, #2 notwendig
\setlength{\gb}{\linewidth}
\addtolength{\gb}{-5.25pt}
\ifthenelse{\equal{\value{itemoddeven}}{0}}{%
\noindent\colorbox{bgeven}{\hskip-3pt\begin{minipage}{\gb}\Qitem[#1]{#2}\end{minipage}}%
\stepcounter{itemoddeven}%
}{%
\noindent\colorbox{bgodd}{\hskip-3pt\begin{minipage}{\gb}\Qitem[#1]{#2}\end{minipage}}%
\setcounter{itemoddeven}{0}%
}
}

\newcommand{\rawcode}[1]{\texttt{#1}}
\newcommand{\code}[2]{\mintinline{#1}{#2}}
\newcommand{\file}[1]{\texttt{\underline{#1}}}

\renewcommand{\theFancyVerbLine}{\tiny \texttt{\textcolor{gray}{\arabic{FancyVerbLine}}}}

\newcommand{\inputcode}[2]{\begin{mdframed}\inputminted[fontsize=\footnotesize,linenos,numbersep=4pt]{#1}{#2}\end{mdframed}}

\newenvironment{solution}{%
    \begin{mdframed}%
        \noindent\textbf{Solu\c{c}\~ao:}\par\medskip
}{%
    \end{mdframed}%
}
\newcommand{\inputsnippet}[4]{\begin{mdframed}[shadow=true,shadowsize=3pt,shadowcolor=gray]\inputminted[fontsize=\footnotesize,linenos,firstline=#2,lastline=#3,numbersep=2pt]{#1}{#4}\end{mdframed}}

\newcommand{\inputsyntax}[2]{\inputminted[fontsize=\footnotesize]{#1}{#2}}



\begin{document}

\begin{center}
	\includegraphics[scale=0.45]{fga}
	% fga.ps: 1087x103 pixel, 72dpi, 38.35x3.63 cm, bb=0 0 1087 103
\end{center}


\vspace{0.1in}

\begin{center}
    \textbf{\huge Paradigmas de Programação}

    \vspace{0.1in}
    \textbf{\Large Programação Vetorial: Lista}
\end{center}\vskip1em

\renewcommand{\Qq}[1]{{#1}}

\Qitem{ \Qq{Analise as funções táticas abaixo, quanto à sua aridade, classificação e semântica, dando
        exemplos em cada caso.  \textbf{Exemplo}: a função tática

            \inputsyntax{apl}{codes/sample.apl}

        em sua forma monádica, é um \textit{fork} que computa a diferença entre o conjugado
        complexo e o vetor unitário na direção de seu argumento:

            \inputsyntax{apl}{codes/samples1.apl}

        Em sua forma diádica é um \textit{fork} que computa a diferença entre a soma e o produto
        de ambos argumentos:

            \inputsyntax{apl}{codes/samples2.apl}

        \textbf{Dica:} utilize o site TryAPL para consultar as versões monádicas e diádicas de
        cada glifo (e também o comando \code{apl}{]help}).}

        \vspace{0.1in}

    \begin{Qlist}
    \item[(a)] \code{apl}{÷*-}

    \begin{solution}
    \textbf{\textit{Fork}} com $f=$ \code{apl}{÷}, $g=$ \code{apl}{*}, $h=$ \code{apl}{-}.
    \begin{itemize}
        \item \textbf{Monádica:} $(\div\omega)^{(-\omega)} = (1/\omega)^{-\omega}=\omega^\omega$.
        \item \textbf{Diádica:} $(\alpha\div\omega)^{(\alpha-\omega)}$.
    \end{itemize}
    \inputsyntax{apl}{codes/sol/vet_a.apl}
    \end{solution}

    \item[(b)] \code{apl}{⍳!+}

    \begin{solution}
    \textbf{\textit{Fork}} com $f=$ \code{apl}{⍳}, $g=$ \code{apl}{!}, $h=$ \code{apl}{+}.
    \begin{itemize}
        \item \textbf{Monádica:} \code{apl}{(⍳⍵)!+⍵}. Para $\omega$ escalar, gera os coeficientes binomiais $\binom{\omega}{1},\ldots,\binom{\omega}{\omega}$ (linha $\omega$ do triângulo de Pascal sem $\binom{\omega}{0}$).
        \item \textbf{Diádica:} \code{apl}{(⍺⍳⍵)!(⍺+⍵)}, calculando $\binom{\alpha+\omega}{k}$, onde $k$ é o índice de $\omega$ no vetor $\alpha$.
    \end{itemize}
    \inputsyntax{apl}{codes/sol/vet_b.apl}
    \end{solution}

    \item[(c)] \code{apl}{×/*}

    \begin{solution}
    \textbf{\textit{Atop}} com $f=$ \code{apl}{×/} e $g=$ \code{apl}{*}.
    \begin{itemize}
        \item \textbf{Monádica:} \code{apl}{×/*⍵}, que calcula $\prod_i e^{\omega_i}=e^{\sum_i\omega_i}$.
        \item \textbf{Diádica:} \code{apl}{×/(⍺*⍵)}, o produto das potências $\prod_i \alpha_i^{\omega_i}$.
    \end{itemize}
    \inputsyntax{apl}{codes/sol/vet_c.apl}
    \end{solution}

    \item[(d)] \code{apl}{⌈÷⌊}

    \begin{solution}
    \textbf{\textit{Fork}} com $f=$ \code{apl}{⌈}, $g=$ \code{apl}{÷}, $h=$ \code{apl}{⌊}.
    \begin{itemize}
        \item \textbf{Monádica:} $\lceil\omega\rceil\div\lfloor\omega\rfloor$ ($1$ para inteiros).
        \item \textbf{Diádica:} $(\alpha\lceil\omega)\div(\alpha\lfloor\omega)$, o maior dividido pelo menor.
    \end{itemize}
    \inputsyntax{apl}{codes/sol/vet_d.apl}
    \end{solution}

    \item[(e)] \code{apl}{+=-}

    \begin{solution}
    \textbf{\textit{Fork}} com $f=$ \code{apl}{+}, $g=$ \code{apl}{=}, $h=$ \code{apl}{-}.
    \begin{itemize}
        \item \textbf{Monádica:} $\overline{\omega}=-\omega$, verdadeiro quando a parte real de $\omega$ é nula.
        \item \textbf{Diádica:} $(\alpha+\omega)=(\alpha-\omega)$, verdadeiro quando $\omega=0$.
    \end{itemize}
    \inputsyntax{apl}{codes/sol/vet_e.apl}
    \end{solution}

    \item[(f)] \code{apl}{×≥÷}

    \begin{solution}
    \textbf{\textit{Fork}} com $f=$ \code{apl}{×}, $g=$ \code{apl}{≥}, $h=$ \code{apl}{÷}.
    \begin{itemize}
        \item \textbf{Monádica:} $\operatorname{sgn}(\omega)\geq 1/\omega$, retornando $1$ para $\omega\geq 1$ ou $-1\leq\omega<0$, e $0$ nos demais casos.
        \item \textbf{Diádica:} $(\alpha\times\omega)\geq(\alpha\div\omega)$.
    \end{itemize}
    \inputsyntax{apl}{codes/sol/vet_f.apl}
    \end{solution}

    \item[(g)] \code{apl}{≡⍴>}

    \begin{solution}
    \textbf{\textit{Fork}} com $f=$ \code{apl}{≡}, $g=$ \code{apl}{⍴}, $h=$ \code{apl}{>}.
    \begin{itemize}
        \item \textbf{Diádica:} \code{apl}{(⍺≡⍵)⍴(⍺>⍵)}. Resulta em \code{apl}{,0} quando os argumentos são idênticos e no vetor vazio caso contrário.
    \end{itemize}
    \inputsyntax{apl}{codes/sol/vet_g.apl}
    \end{solution}

    \item[(h)] \code{apl}{!+○}

    \begin{solution}
    \textbf{\textit{Fork}} com $f=$ \code{apl}{!}, $g=$ \code{apl}{+}, $h=$ \code{apl}{○}.
    \begin{itemize}
        \item \textbf{Monádica:} $\omega!+\pi\omega$.
        \item \textbf{Diádica:} $\binom{\omega}{\alpha}+(\alpha\,\texttt{○}\,\omega)$, em que \code{apl}{⍺○⍵} aplica a função circular selecionada por $\alpha$.
    \end{itemize}
    \inputsyntax{apl}{codes/sol/vet_h.apl}
    \end{solution}

    \item[(i)] \code{apl}{∨-∧}

    \begin{solution}
    \textbf{\textit{Fork}} com $f=$ \code{apl}{∨}, $g=$ \code{apl}{-}, $h=$ \code{apl}{∧}.
    \begin{itemize}
        \item \textbf{Diádica:} $(\alpha\lor\omega)-(\alpha\land\omega)$. Calcula $\operatorname{mdc}(\alpha,\omega)-\operatorname{mmc}(\alpha,\omega)$ para inteiros positivos e o ou exclusivo para booleanos.
    \end{itemize}
    \inputsyntax{apl}{codes/sol/vet_i.apl}
    \end{solution}

    \item[(j)] \code{apl}{≤≠≥}

    \begin{solution}
    \textbf{\textit{Fork}} com $f=$ \code{apl}{≤}, $g=$ \code{apl}{≠}, $h=$ \code{apl}{≥}.
    \begin{itemize}
        \item \textbf{Diádica:} $(\alpha\leq\omega)\neq(\alpha\geq\omega)$, equivalente a $\alpha\neq\omega$.
    \end{itemize}
    \inputsyntax{apl}{codes/sol/vet_j.apl}
    \end{solution}

    \item[(k)] \code{apl}{⍋⍒}

    \begin{solution}
    \textbf{\textit{Atop}} com $f=$ \code{apl}{⍋} e $g=$ \code{apl}{⍒}.
    \begin{itemize}
        \item \textbf{Monádica:} \code{apl}{⍋(⍒⍵)}, que retorna a posição de cada elemento de $\omega$ na ordenação decrescente (inversa da permutação \code{apl}{⍒⍵}).
        \item \textbf{Diádica:} \code{apl}{⍋(⍺⍒⍵)}, que faz o mesmo usando o vetor de caracteres $\alpha$ como ordem de comparação.
    \end{itemize}
    \inputsyntax{apl}{codes/sol/vet_k.apl}
    \end{solution}

    \item[(l)] \code{apl}{⍴⍴≡}

    \begin{solution}
    \textbf{\textit{Fork}} com $f=$ \code{apl}{⍴}, $g=$ \code{apl}{⍴}, $h=$ \code{apl}{≡}.
    \begin{itemize}
        \item \textbf{Monádica:} \code{apl}{(⍴⍵)⍴(≡⍵)}, um \emph{array} com a forma de $\omega$ preenchido com sua profundidade.
        \item \textbf{Diádica:} \code{apl}{(⍺⍴⍵)⍴(⍺≡⍵)}, um \emph{array} de forma \code{apl}{⍺⍴⍵} preenchido com $1$ se os argumentos forem idênticos e $0$ caso contrário.
    \end{itemize}
    \inputsyntax{apl}{codes/sol/vet_l.apl}
    \end{solution}

    \item[(m)] \code{apl}{⊃⊂⍳}

    \begin{solution}
    \textbf{\textit{Fork}} com $f=$ \code{apl}{⊃}, $g=$ \code{apl}{⊂}, $h=$ \code{apl}{⍳}.
    \begin{itemize}
        \item \textbf{Monádica:} \code{apl}{(⊃⍵)⊂(⍳⍵)}, particionando \code{apl}{⍳⍵} com \code{apl}{⊃⍵} partições por posição.
        \item \textbf{Diádica:} \code{apl}{(⍺⊃⍵)⊂(⍺⍳⍵)}, particionando os índices de $\omega$ em $\alpha$ conforme o valor selecionado por \code{apl}{⍺⊃⍵}.
    \end{itemize}
    \inputsyntax{apl}{codes/sol/vet_m.apl}
    \end{solution}

    \item[(n)] \inputsyntax{apl}{codes/scan.apl}

    \begin{solution}
    \textbf{\textit{Atop}} com $f=$ \code{apl}{⍟} e $g=$ \texttt{×\textbackslash}.
    \begin{itemize}
        \item \textbf{Monádica:} \texttt{⍟(×\textbackslash⍵)}, o logaritmo natural dos produtos acumulados, equivalente à soma acumulada dos logaritmos (\texttt{+\textbackslash⍟⍵}).
    \end{itemize}
    \inputsyntax{apl}{codes/sol/vet_n.apl}
    \end{solution}

    \end{Qlist}
}


\end{document}
